problem:  
Let \( (X, J, \omega, \Omega, H) \) be a compact complex threefold satisfying the **Hull–Strominger system**:
\[
\begin{cases}
H = i(\bar{\partial} - \partial)\omega, \\[4pt]
d(\|\Omega\|_\omega\, \omega^2) = 0, \\[4pt]
dH = \dfrac{\alpha'}{4}\big(\mathrm{Tr}(R \wedge R) - \mathrm{Tr}(F \wedge F)\big),
\end{cases}
\]
where \(R\) is the Chern curvature of the Hermitian metric on \(T^{1,0}X\), and \(F\) is the curvature of a Hermitian–Yang–Mills connection on a holomorphic bundle \(E \to X\).

Let \(\widehat{H} \in \widehat{H}^3(X,\mathbb{Z})\) denote the **differential cohomology class** whose curvature is the 3‑form \(H\), satisfying the Bianchi identity above.  
The quantization condition requires \(\widehat{H}\) to project to an **integral cohomology class** \([H] \in H^3(X,\mathbb{Z})\), which may have a **nontrivial torsion component**.

Deformations of the Hull–Strominger system are known to be controlled by a differential graded complex \(\mathcal{Q}_{X,H}^\bullet\) (as constructed by de la Ossa–Svanes). When flux quantization is ignored, unobstructedness of deformations is governed by classes in \(H^2(\mathcal{Q}_{X,H}^\bullet)\).  
However, enforcing the flux constraint in differential cohomology restricts admissible deformations to those \((\delta J, \delta \omega, \delta H)\) whose curvature 3‑form variation satisfies
\[
[\widehat{H + \delta H}] = [\widehat{H}] \ \ \text{in } \widehat{H}^3(X,\mathbb{Z}),
\]
so that the **integral and torsion components remain fixed** under deformation.

> **Question:**  
> When this differential‑cohomological constraint is imposed, can a nontrivial torsion component in \([H] \in H^3(X,\mathbb{Z})_{\mathrm{tors}}\) give rise to *new analytic obstruction classes* in \(H^2(\mathcal{Q}_{X,H}^\bullet)\) that would otherwise vanish?  
> Equivalently, does discrete torsion in the quantized flux modify the unobstructedness of the analytic deformation complex for the Strominger system, once one enforces preservation of the differential cohomology class \(\widehat{H}\)?

Why is it a "good" problem:  
This question is both precise and profound. It targets a fundamental but unresolved interface between **analytic deformation theory** and **integral topology** in heterotic compactifications.  
The deformation complex \(\mathcal{Q}_{X,H}^\bullet\) is analytic, defined over smooth forms, while the flux quantization condition lives in **differential cohomology**, capturing discrete torsion information invisible to de Rham cohomology. Determining whether torsion data can re‑enter the analytic obstruction theory once flux constraints are imposed would illuminate how **discrete topological fluxes** interact with **smooth moduli spaces** in non‑Kähler Calabi–Yau geometry.  
A positive answer would reveal a new coupling between **differential topology** and **analysis**, influencing not only geometric flows and moduli stability but also the mathematical foundations of flux compactification.  
Its expression is compact, logically sound, and conceptually unifying—qualities that make it an exemplary "good" research problem.